\uuid{R1mK}
\titre{Vecteurs}
\auteur{bodin ; hanani ; mzari}
\organisation{exo7}
\licence{CC BY-NC-SA 4.0}
\datecreate{2020-01-20}
\niveau{L1}
\module{Géométrie}
\chapitre{Géométrie du plan}
\sousChapitre{Distance, norme, produit scalaire}
\difficulte{1}
\theme{vecteurs}
\exo7qcmid{lille-1.geometrie-plane.008}

\contenu{
\begin{question}
\qclassification{140.01,140.02}
\qtags{motcle=vecteurs}
Soit $D$ la droite passant par l'origine et par le point $A(1,1)$. Quelles sont les bonnes réponses ?
\begin{answers}
    \good{$\vec{u}(1,1)$ est un vecteur directeur de $D$.}
    \bad{$\vec{u}(1,1)$ est un vecteur normal à $D$.}
    \good{$y=x$ est une équation cartésienne de $D$.}
    \bad{$x+y=0$ est une équation cartésienne de $D$.}
\end{answers}
\begin{explanations}
La droite $D$ est dirigée par le vecteur $\overrightarrow{OA}=\vec{u}(1,1)$ et $M(x,y)\in D \Leftrightarrow \mbox{det}(\overrightarrow{OM},\overrightarrow{OA})=0\Leftrightarrow x-y=0$. Ceci donne une équation cartésienne de $D$.
\end{explanations}
\end{question}
}
